\subsection{Market Sample and Public Data Design}

The empirical design is built around one principle: every variable used in the forecasting and warning system must be publicly accessible, chronologically alignable, and reproducible without proprietary terminals or vendor-specific feeds. This requirement is not cosmetic. In volatility forecasting, reported gains can be materially affected by hidden differences in data timing, revision handling, and sample construction. A public-data design therefore sharpens both the scientific and practical value of the study.

The market sample contains the S\&P~500, Nasdaq-100, and Dow Jones Industrial Average. These three indices provide a compact but informative cross-section for the problem studied here. The S\&P~500 serves as the broad U.S.\ market benchmark; the Nasdaq-100 represents a technology- and growth-concentrated segment with distinct volatility dynamics; and the Dow Jones Industrial Average provides a price-weighted blue-chip benchmark. Together they form a natural setting for testing both within-index multiscale dynamics and cross-index spillover. Daily OHLC prices are collected from Stooq. Public macro-financial predictors, implied-volatility proxies, rates, term spreads, financial conditions, and recession indicators are collected from FRED. ETF-based trading-volume proxies are drawn from SPY, QQQ, and DIA. The raw-data snapshot used by the experimental configuration is fixed to 17 March 2026 so that the reported results correspond to a stable and reproducible input state rather than a moving download target.

\begin{table}[H]
\caption{Public data blocks used in the baseline and robustness specifications. The table is organized by predictor block, source, and empirical role in the forecasting pipeline.}
\label{tab:data_sources}
\centering
\small
\setlength{\tabcolsep}{4pt}
\renewcommand{\arraystretch}{1.15}
\begin{tabularx}{\textwidth}{>{\raggedright\arraybackslash}p{0.16\textwidth}>{\raggedright\arraybackslash}X>{\raggedright\arraybackslash}p{0.16\textwidth}>{\raggedright\arraybackslash}p{0.24\textwidth}}
\toprule
Block & Variables / symbols & Source & Main role \\
\midrule
Equity indices & S\&P 500, Nasdaq-100, DJIA & Stooq & Daily OHLC prices and volatility targets \\
Liquidity proxies & SPY, QQQ, DIA & Stooq & ETF-based volume and trading-activity proxies \\
Implied volatility & VIXCLS, VXNCLS, VXDCLS & FRED (CBOE series) & Index-specific forward-looking risk proxies \\
Rates and conditions & DFF, DGS10, T10Y3M, NFCI, USRECD & FRED & Macro-financial and regime controls \\
Expanded public risk data & DGS2, BAMLH0A0HYM2, BAMLC0A0CM, USEPUINDXD, RVXCLS, VXVCLS & FRED & Credit, policy-uncertainty, and volatility-term-structure robustness block \\
\bottomrule
\end{tabularx}
\end{table}


The primary sample spans 25 February 2005 to 31 December 2025. This starting point balances three requirements: long enough history for rolling estimation and multiscale decomposition, comparable availability across the three indices, and consistent access to ETF-based liquidity proxies. The out-of-sample forecasting period begins on 4 January 2016, leaving a sufficiently long pre-2016 history for training and feature construction. All predictors are aligned to index trading days, and non-trading days are excluded from the panel. Daily and weekly macro-financial series are forward-filled only after they become historically observable. Monthly variables with ambiguous release timing are excluded from the main benchmark specification to avoid avoidable timing disputes in a causal design.

\subsection{Volatility Target Construction}

The target of interest is future average stock-index volatility constructed from daily OHLC data. The main specification uses the Rogers--Satchell estimator because it retains range information while remaining more robust than close-to-close volatility under nonzero drift \cite{rogers1991estimating}. For each forecast horizon \(h\in\{1,5,10\}\), the target is defined as the average volatility proxy over the next \(h\) trading days. This choice aligns the regression task with forward-looking risk monitoring rather than same-day measurement.

The target is modeled on the \(\log(1+y)\) scale with a small positive constant for numerical stability and then transformed back to the original scale for evaluation. This preserves positivity while limiting the leverage of exceptionally large observations. The main target is complemented by a robustness specification based on the Parkinson estimator \cite{parkinson1980extreme}. Using both Rogers--Satchell and Parkinson is substantively useful. If the proposed framework improves only under one narrowly chosen volatility proxy, its contribution would be weak. If its relative advantage persists when the target emphasizes high--low range information differently, that is stronger evidence that the gain comes from representation quality rather than from target-specific tuning. The classical OHLC estimators of Parkinson, Garman--Klass, Rogers--Satchell, and Yang--Zhang provide the statistical background for this design \cite{parkinson1980extreme,garman1980estimation,rogers1991estimating,yang2000drift,andersen2003rvol}.

\subsection{Predictor Architecture}

The predictor set is organized into four blocks so that the empirical design can separate persistence, raw market information, public exogenous risk information, and multiscale representation.

The first block contains persistence features. These include the current volatility proxy and HAR-style aggregates over the last 1, 5, and 22 trading days. This block is essential because short-run volatility forecasting is dominated by persistence, and any richer model must prove that it adds value beyond this structure rather than simply re-encoding it.

The second block contains raw market-transition and technical signals. These include close-to-close returns, intraday returns, overnight gaps, RSI, MACD, ATR, and short-window stress summaries computed over 2-, 3-, and 5-day windows. This block is designed to capture local market transitions that may not be fully represented by the volatility target alone.

The third block contains public exogenous risk variables. These include implied-volatility series, interest-rate and spread variables, financial-conditions indicators, and a recession indicator. In the expanded-data robustness setting, this block is enlarged with credit-spread, policy-uncertainty, and volatility-term-structure variables to test whether the proposed framework benefits from a richer risk-information environment rather than only from a carefully limited predictor set.

The fourth block contains causal multiscale wavelet summaries. These are extracted from the target volatility series, absolute returns, and the index-specific implied-volatility proxy. This block is where the CMWSL framework departs most clearly from standard tabular forecasting. Instead of treating raw lags as the only temporal representation, it summarizes the recent scale-wise behaviour of volatility-related signals through causal wavelet coefficients and energy statistics. The modular block design is deliberate: it allows the empirical analysis to distinguish baseline persistence, conventional public information, and the incremental value of multiscale representation.

Two additional design choices deserve emphasis. First, the main specification uses only variables whose timing can be handled cleanly in a daily walk-forward setting; this reduces the scope for hidden look-ahead effects. Second, the robustness extensions are intentionally challenging: a broader public-information panel and an alternative OHLC target both raise the standard for claiming that multiscale learning provides genuine incremental value.
